\documentclass[conference]{IEEEtran}
\IEEEoverridecommandlockouts
% The preceding line is only needed to identify funding in the first footnote. If that is unneeded, please comment it out.
\usepackage{cite}
\usepackage{amsmath,amssymb,amsfonts}
\usepackage{algorithmic}
\usepackage{graphicx}
\usepackage{textcomp}
\usepackage{xcolor}
\usepackage{setspace}
\usepackage{float}
\usepackage{empheq}

\newsavebox{\equationBox}
\newlength{\equationBoxCeiling}
\newlength{\equationBoxFloor}
\NewDocumentCommand{\boxThisEquation}{O{0pt} O{0pt} O{0pt} O{0pt} m}{%
    \sbox{\equationBox}{#5}%
    \setlength{\equationBoxCeiling}{\ht\equationBox}%
    \addtolength{\equationBoxCeiling}{#1}%
    \setlength{\ht\equationBox}{\equationBoxCeiling}%
    %
    \setlength{\equationBoxFloor}{\dp\equationBox}%
    \addtolength{\equationBoxFloor}{#2}%
    \setlength{\dp\equationBox}{\equationBoxFloor}%
    %
    \fbox{\hspace{2mm+#3}\usebox{\equationBox}\hspace{2mm+#4}}%
}

\def\BibTeX{{\rm B\kern-.05em{\sc i\kern-.025em b}\kern-.08em
    T\kern-.1667em\lower.7ex\hbox{E}\kern-.125emX}}
\begin{document}

\title{Environment Similarity Index: Quantifying map-environment divergences for robot localization performance estimation and improvement\\
\thanks{This work has been funded by the Danish Innovation Fund.}
}

\author{\IEEEauthorblockN{Christian Rosendahl Dam}
\IEEEauthorblockA{\textit{University of Southern Denmark} \\
\textit{The Maersk Mc-Kinney Moller Institute}\\
Odense, Denmark \\
crda@mir-robots.com}
\and
\IEEEauthorblockN{Daniel Suvei}
\IEEEauthorblockA{\textit{Traffic and Navigation Team} \\
\textit{Teradyne Robotics}\\
Odense, Denmark \\
dsu@mir-robots.com}
\and
\IEEEauthorblockN{Leon Bodenhagen}
\IEEEauthorblockA{\textit{University of Southern Denmark} \\
\textit{The Maersk Mc-Kinney Moller Institute}\\
Odense, Denmark \\
lebo@mmmi.sdu.dk}
}

\maketitle

\begin{abstract}
Should take about this much space: \newline
abc abc abc abc abc abc abc abc abc abc abc abc abc abc abc abc abc abc abc abc abc abc abc abc abc abc abc abc abc abc abc abc abc abc abc abc abc abc abc abc abc abc abc abc
abc abc abc abc abc abc abc abc abc abc abc abc abc abc abc abc abc abc abc abc abc abc abc abc abc abc abc abc abc abc abc abc abc abc abc abc abc abc abc abc abc abc abc abc
abc abc abc abc abc abc abc abc abc abc abc abc abc abc abc abc abc abc abc abc abc abc abc abc abc abc abc abc abc abc abc abc abc abc abc abc abc abc abc abc abc abc abc abc
abc abc abc abc abc abc abc abc abc abc abc abc abc abc abc abc abc abc abc abc abc abc abc abc abc abc abc abc abc abc abc abc abc abc abc abc abc abc abc abc abc abc abc abc. 
\newline
\end{abstract}

\begin{IEEEkeywords}
Localization, Mapping, AMR
\end{IEEEkeywords}

\section{Introduction}
Robust, lifelong localization is a fundamental ability for all Autonomous Mobile Robots (AMRs). For many years, the gold standard to achieve this has been to use some variation 
of Monte Carlo Localization, as it strikes a good balance between cost and performance. It can perform global localization, and is multi-modal. For indoor industrial applications, 
it works well in many cases. It does however, have a major drawback: It assumes the environment is static. As soon as it' surroundings have changed too much, the particle filter 
diverges and fails, because it does not update it's representation. Making a system that is robust to change is still an ongoing challenge (ref). Some would argue that 
Simultaneous Localization and Mapping (SLAM) is the best way to solve this; but SLAM suffers in reliability. No system is perfect, and sooner or later, an error will find its way 
into the system. The resulting erroneous pose will propagate into the resulting map, reducing accuracy or cause a complete loss of localization. Since there is no ground truth 
to inform the system when the errors was introduced, it is impossible to recover the last valid map state. Therefore, this paper suggests a different approach; instead of updating 
the map, the environment is analysed to detect changes between what is seen and what is expected. This info is then used to alter the behavior of the localization system 
- in this case, adaptive MCL (AMCL). To do this, we need a metric that defines how similar (or different) the observed environment is from its representation (the map). 
However; to the best of our knowledge, no such metric exists. We hypothesize that such a metric can be used to infer the performance of any given localization module that takes 
the static world assumption, and that it can be used online to adjust its behaviour, thereby reducing the pose error. Furthermore, it can be used to generate a heatmap of 
an environment, which can then be interpreted by a human operator, in order to identify local regions of high divergence. These are the contributions of this paper:

\subsection{Contribution}
This is the main contribution of this paper. We propose the Environment Similarity Index (ESI): An index that quantifies how well the currently observed environment matches the map.
We also show how the ESI correlates to changes in the environment, both on simulated and real-world data, and how these changes impacts AMCL's performance. 
In addition, we show how the ESI can be used by a user to easily identify areas in their site that have changed since mapping, and therefore may increase localiztion errors. 
Finally, we show how the ESI can be used to improve AMCL's performance by using it to dynamically tune AMCL's parameters.


\section{Related Work}
- Bayes filtering / Probabilistic map update
- Scan-BIM divergences
- Occupancy grid vs feature maps
- 





\begin{table}[htbp]
\caption{Table Type Styles}
\begin{center}
\begin{tabular}{|c|c|c|c|}
\hline
\textbf{Table}&\multicolumn{3}{|c|}{\textbf{Table Column Head}} \\
\cline{2-4} 
\textbf{Head} & \textbf{\textit{Table column subhead}}& \textbf{\textit{Subhead}}& \textbf{\textit{Subhead}} \\
\hline
copy& More table copy$^{\mathrm{a}}$& &  \\
\hline
\multicolumn{4}{l}{$^{\mathrm{a}}$Sample of a Table footnote.}
\end{tabular}
\label{tab1}
\end{center}
\end{table}


\section*{Acknowledgment}
...

\newpage
\begin{thebibliography}{00}
\bibitem{b1} M. Young, The Technical Writer's Handbook. Mill Valley, CA: University Science, 1989.
\bibitem{b2} J. Doe, A. Smith, and B. Johnson, ``Lorem ipsum dolor sit amet, consectetur adipiscing elit,'' in Proc. IEEE Int. Conf. Robotics Automat., 2020, pp. 123--128.
\bibitem{b3} C. Williams, D. Brown, and E. Davis, ``Sed do eiusmod tempor incididunt ut labore,'' IEEE Trans. Robot., vol. 35, no. 4, pp. 456--467, Aug. 2019.
\bibitem{b4} F. Miller, G. Wilson, and H. Taylor, ``Ut enim ad minim veniam, quis nostrud exercitation,'' in Proc. IEEE/RSJ Int. Conf. Intell. Robots Syst., 2021, pp. 234--241.
\bibitem{b5} I. Anderson, J. Martinez, and K. Lee, ``Ullamco laboris nisi ut aliquip ex ea commodo consequat,'' IEEE Robot. Automat. Lett., vol. 6, no. 2, pp. 789--795, Apr. 2021.
\bibitem{b6} L. Garcia, M. Rodriguez, and N. Patel, ``Duis aute irure dolor in reprehenderit,'' in Proc. IEEE Int. Conf. Robot. Automat., 2018, pp. 567--574.
\bibitem{b7} O. Thompson, P. White, and Q. Harris, ``In voluptate velit esse cillum dolore eu fugiat,'' IEEE Trans. Automat. Sci. Eng., vol. 17, no. 3, pp. 1234--1245, Jul. 2020.
\bibitem{b8} R. Clark, S. Lewis, and T. Walker, ``Nulla pariatur excepteur sint occaecat cupidatat,'' in Proc. IEEE Conf. Decision Control, 2019, pp. 345--352.
\bibitem{b9} U. Young, V. Hall, and W. Allen, ``Non proident, sunt in culpa qui officia deserunt,'' IEEE Trans. Syst., Man, Cybern., Syst., vol. 50, no. 8, pp. 2345--2356, Aug. 2020.
\bibitem{b10} X. King, Y. Wright, and Z. Scott, ``Mollit anim id est laborum,'' in Proc. IEEE Int. Conf. Robotics Automat., 2022, pp. 678--685.
\bibitem{b11} A. Adams, B. Baker, and C. Carter, ``Sed ut perspiciatis unde omnis iste natus error,'' IEEE Robot. Automat. Mag., vol. 27, no. 2, pp. 89--97, Jun. 2020.
\bibitem{b12} D. Evans, E. Foster, and F. Green, ``Sit voluptatem accusantium doloremque laudantium,'' in Proc. IEEE/RSJ Int. Conf. Intell. Robots Syst., 2020, pp. 456--463.
\bibitem{b13} G. Hill, H. Jackson, and I. Jones, ``Totam rem aperiam, eaque ipsa quae ab illo inventore,'' IEEE Trans. Robot., vol. 36, no. 5, pp. 1234--1245, Oct. 2020.
\bibitem{b14} J. Kelly, K. Lopez, and L. Moore, ``Veritatis et quasi architecto beatae vitae dicta sunt,'' in Proc. IEEE Int. Conf. Robot. Automat., 2021, pp. 789--796.
\bibitem{b15} M. Nelson, N. Owens, and O. Parker, ``Explicabo nemo enim ipsam voluptatem quia voluptas,'' IEEE Robot. Automat. Lett., vol. 5, no. 4, pp. 2345--2352, Oct. 2020.
\bibitem{b16} P. Quinn, R. Reed, and S. Ross, ``Sit aspernatur aut odit aut fugit, sed quia consequuntur,'' in Proc. IEEE Conf. Decision Control, 2020, pp. 567--574.
\bibitem{b17} T. Stewart, U. Turner, and V. Underwood, ``Magni dolores eos qui ratione voluptatem sequi nesciunt,'' IEEE Trans. Automat. Sci. Eng., vol. 18, no. 2, pp. 3456--3467, Apr. 2021.
\bibitem{b18} W. Vance, X. Vega, and Y. Vincent, ``Neque porro quisquam est, qui dolorem ipsum quia dolor sit amet,'' in Proc. IEEE Int. Conf. Robotics Automat., 2019, pp. 890--897.
\bibitem{b19} Z. Ward, A. Watson, and B. Webb, ``Consectetur, adipisci velit, sed quia non numquam eius modi tempora,'' IEEE Trans. Syst., Man, Cybern., Syst., vol. 51, no. 6, pp. 4567--4578, Jun. 2021.
\bibitem{b20} C. West, D. Wood, and E. Wright, ``Incidunt ut labore et dolore magnam aliquam quaerat voluptatem,'' in Proc. IEEE/RSJ Int. Conf. Intell. Robots Syst., 2021, pp. 123--130.
\bibitem{b21} F. Yang, G. Young, and H. Zhang, ``Ut enim ad minima veniam, quis nostrum exercitationem ullam corporis,'' IEEE Robot. Automat. Mag., vol. 28, no. 3, pp. 123--134, Sep. 2021.
\bibitem{b22} I. Zhao, J. Zhou, and K. Zhu, ``Suscipit laboriosam, nisi ut aliquid ex ea commodi consequatur,'' in Proc. IEEE Int. Conf. Robot. Automat., 2020, pp. 234--241.
\bibitem{b23} L. Abbott, M. Bell, and N. Cooper, ``Quis autem vel eum iure reprehenderit qui in ea voluptate velit esse,'' IEEE Trans. Robot., vol. 37, no. 1, pp. 5678--5689, Feb. 2021.
\bibitem{b24} O. Davis, P. Edwards, and Q. Ellis, ``Quam nihil molestiae consequatur, vel illum qui dolorem eum fugiat,'' in Proc. IEEE Conf. Decision Control, 2021, pp. 345--352.
\bibitem{b25} R. Fisher, S. Ford, and T. Fox, ``Quo voluptas nulla pariatur at vero eos et accusamus,'' IEEE Robot. Automat. Lett., vol. 6, no. 3, pp. 6789--6796, Jul. 2021.
\bibitem{b26} U. Grant, V. Gray, and W. Green, ``Et iusto odio dignissimos ducimus qui blanditiis praesentium,'' in Proc. IEEE Int. Conf. Robotics Automat., 2021, pp. 456--463.
\bibitem{b27} X. Hayes, Y. Hughes, and Z. Hunt, ``Voluptatum deleniti atque corrupti quos dolores et quas molestias,'' IEEE Trans. Automat. Sci. Eng., vol. 19, no. 1, pp. 7890--7901, Jan. 2022.
\bibitem{b28} A. Ingram, B. Irwin, and C. Irwin, ``Excepturi sint occaecati cupiditate non provident, similique sunt in culpa,'' in Proc. IEEE/RSJ Int. Conf. Intell. Robots Syst., 2022, pp. 567--574.
\bibitem{b29} D. Jenkins, E. Jensen, and F. Johnson, ``Qui officia deserunt mollitia animi, id est laborum et dolorum fuga,'' IEEE Trans. Syst., Man, Cybern., Syst., vol. 52, no. 4, pp. 2345--2356, Apr. 2022.
\bibitem{b30} G. Klein, H. Knight, and I. Knox, ``Et harum quidem rerum facilis est et expedita distinctio,'' in Proc. IEEE Int. Conf. Robot. Automat., 2022, pp. 678--685.
\end{thebibliography}
\vspace{12pt}

\end{document}